\documentclass[a4paper,12pt]{article}
\usepackage[T2A]{fontenc}
\usepackage[utf8]{inputenc}
\usepackage[russian]{babel}
\usepackage{cmap}
\usepackage{amsmath,amssymb,mathtools}
\usepackage{geometry}
\geometry{left=21mm,right=21mm,top=20mm,bottom=22mm,headheight=16pt,headsep=9mm}
\usepackage{microtype}
\usepackage{setspace}
\setstretch{1.15}
\usepackage{enumitem}
\usepackage{tabularx,booktabs,array}
\usepackage{tikz}
\usetikzlibrary{calc}
\usepackage{xcolor}
\usepackage{fancyhdr}
\usepackage{hyperref}
\usepackage{needspace}
\definecolor{accent}{HTML}{00848A}
\definecolor{darkaccent}{HTML}{16585E}
\definecolor{textgray}{HTML}{3E4143}
\definecolor{linegray}{HTML}{D8DEDE}
\hypersetup{colorlinks=true,linkcolor=darkaccent,pdfauthor={Лёвин Артём Александрович},pdftitle={Комплексное повторение. 10 упражнений — Анастасия Павлова}}
\color{textgray}
\setlength{\parindent}{0pt}
\setlength{\parskip}{5pt}
\pagestyle{fancy}
\fancyhf{}
\fancyhead[L]{\small Комплексное повторение}
\fancyhead[R]{\small Анастасия Павлова}
\fancyfoot[C]{\small\thepage}
\renewcommand{\headrulewidth}{0.3pt}
\renewcommand{\footrulewidth}{0pt}
\newcommand{\sectionline}[1]{%
  \par\medskip\Needspace{4\baselineskip}
  {\Large\color{darkaccent}#1\par}\vspace{2mm}
  {\color{accent}\rule{\linewidth}{0.55pt}}\par\vspace{2mm}}
\newcommand{\exercise}[3]{%
  \Needspace{6\baselineskip}
  {\large\bfseries\color{darkaccent}#1.\quad #2\par}\vspace{1.5mm}
  #3\par}
\newcommand{\worklines}[1]{%
  \vspace{5mm}
  \begingroup\color{linegray}%
  \foreach \x in {1,...,#1}{\par\vspace{8mm}\rule{\linewidth}{0.25pt}}%
  \endgroup
}
\newcommand{\pairpage}[1]{\clearpage\sectionline{#1}}
\begin{document}
\begin{titlepage}
\thispagestyle{empty}
\centering
\vspace*{24mm}
{\large\color{darkaccent}МАТЕМАТИКА\par}
\vspace{19mm}
{\LARGE Сборник упражнений\par}
\vspace{12mm}
{\LARGE\color{accent}Комплексное повторение\par}
\vspace{5mm}
{\large Степени, функции, задачи ЕГЭ, финансы и вероятность\par}
\vspace{14mm}
{\normalsize Подготовка к ЕГЭ по профильной математике\par}
\vfill
{\large Анастасия Павлова\par}
\vspace{10mm}
{\normalsize 16 ноября 2026 г.\par}
\vspace{20mm}
{\small Лёвин Артём Александрович\par}
{\small эксклюзивно для Анастасии Павловой\par}
\vspace{16mm}
\begin{tikzpicture}
  \draw[accent,line width=0.9pt]
  (-5.1,0) .. controls (-3.5,0.70) and (-2,-0.55) .. (0,0)
  .. controls (1.7,0.65) and (3.5,-0.45) .. (5.1,0);
\end{tikzpicture}
\end{titlepage}

\sectionline{Маршрут повторения}
В сборнике 10 самостоятельных упражнений. Каждое проверяет отдельный навык, который использовался на занятиях. Задачи расположены от базовых преобразований к комплексным моделям.

\vspace{4mm}
\renewcommand{\arraystretch}{1.5}
\begin{tabularx}{\linewidth}{@{}>{\bfseries}p{8mm}X p{36mm}@{}}
\toprule
№ & Тема & Период занятий \\
\midrule
1 & Степени и показательные уравнения & Август — сентябрь \\
2 & Формулы и пересечения графиков & Сентябрь \\
3 & Векторы на координатной плоскости & Сентябрь \\
4 & Движение по воде & Сентябрь \\
5 & Совместная работа & Октябрь \\
6 & Смеси, растворы, испарение & Октябрь \\
7 & Последовательные проценты & Октябрь \\
8 & Вклад с капитализацией и пополнением & Октябрь — ноябрь \\
9 & Сравнение схем погашения кредита & Октябрь — ноябрь \\
10 & Сочетания и вероятность & Ноябрь \\
\bottomrule
\end{tabularx}

\vspace{8mm}
\sectionline{Как работать с заданиями}
\begin{enumerate}[label=\textcolor{accent}{\arabic*.},leftmargin=8mm,itemsep=4pt]
\item Записывайте ключевую математическую модель: формулу, систему или уравнение.
\item Проверяйте ограничения: знаменатели, допустимые скорости и времена, ставки, смысл полученных корней.
\item В задачах с вероятностью сначала определяйте, различается ли порядок выбора.
\item Сверяйте итоговый ответ с таблицей на последней странице только после самостоятельного решения.
\end{enumerate}

\vfill
{\small\color{darkaccent}Рекомендация: выполните все десять упражнений за один подход, затем вернитесь к задачам, в которых возникли затруднения.}

\pairpage{Алгебра и функции}
\exercise{1}{Степени и показательное уравнение}{
Решите уравнение
\[
9^{x-1}=\sqrt{27}\cdot3^{x-2}.
\]
Запишите переход к одному основанию и найдите $x$.
}
\worklines{4}
\vspace{12mm}
\exercise{2}{Функция по двум точкам}{
Линейная функция $f(x)=kx+b$ проходит через точки $A(-2;7)$ и $B(2;-1)$.
Найдите коэффициенты $k$, $b$, значение $f(-3)$ и координаты точки пересечения её графика с прямой $y=x+3$.
}
\worklines{4}

\pairpage{Координаты и движение}
\exercise{3}{Векторы и угол}{
На координатной плоскости заданы точки
\[
A(-2;1),\quad B(4;9),\quad C(6;-5).
\]
Найдите координаты векторов $\overrightarrow{AB}$ и $\overrightarrow{AC}$, длину $\overrightarrow{AB}$ и величину угла $BAC$.
}
\worklines{4}
\vspace{12mm}
\exercise{4}{Движение катера}{
Катер проходит $24$ км против течения на $2$ часа дольше, чем те же $24$ км по течению. Скорость течения равна $3$ км/ч. Найдите собственную скорость катера. Укажите ограничение на неизвестную скорость.
}
\worklines{4}

\pairpage{Совместная работа и смеси}
\exercise{5}{Два исполнителя}{
Первый мастер может выполнить заказ на $6$ часов быстрее второго. Работая вместе с постоянной производительностью, они выполняют заказ за $4$ часа. За сколько часов каждый мастер выполнит заказ самостоятельно?
}
\worklines{4}
\vspace{12mm}
\exercise{6}{Испарение и смешивание}{
В сосуде было $8$ кг $30\%$-го раствора соли. Из него испарили $2$ кг воды, затем добавили $4$ кг $10\%$-го раствора той же соли. Найдите массовую долю соли в полученной смеси (в процентах).
}
\worklines{4}

\pairpage{Проценты и банковские вклады}
\exercise{7}{Два изменения цены}{
Цену товара сначала увеличили на $20\%$, затем полученную цену уменьшили на $10\%$. После этих изменений товар стал стоить $5400$ руб. Найдите первоначальную цену и суммарное изменение цены в процентах относительно первоначальной.
}
\worklines{4}
\vspace{12mm}
\exercise{8}{Вклад с неизвестной ставкой}{
В начале первого года на вклад внесли $600$ тыс. руб. В конце первого года после начисления процентов вклад пополнили на $120$ тыс. руб. В конце второго года на счёте оказалось $858$ тыс. руб. Ставка оба года была одинаковой, проценты капитализировались один раз в конце года. Найдите годовую процентную ставку.
}
\worklines{4}

\pairpage{Финансы, комбинаторика и вероятность}
\exercise{9}{Две схемы погашения кредита}{
Кредит $110$ тыс. руб. выдан на $2$ года под $20\%$ годовых. Проценты начисляются на остаток долга перед очередным ежегодным платежом.

Сравните два варианта полного погашения: (а) два равных ежегодных платежа; (б) ежегодно погашается половина первоначального долга плюс проценты на остаток. Найдите платежи и переплату для обоих вариантов; определите, какая схема дешевле и на сколько.
}
\worklines{4}
\vspace{8mm}
\exercise{10}{Выбор групп и вероятность}{
Шесть участников случайным образом распределяют на две \emph{неименованные} группы по три человека; все такие разбиения равновероятны. Сколько существует разбиений? Какова вероятность того, что два определённых участника, Аня и Борис, окажутся в одной группе?
}
\worklines{4}

\clearpage
\sectionline{Ответы для самопроверки}
В таблице приведены итоговые результаты и короткая контрольная идея. Полные решения намеренно оставлены для самостоятельной работы.\par
\vspace{3mm}
\renewcommand{\arraystretch}{1.58}
\small
\begin{tabularx}{\linewidth}{@{}>{\bfseries}p{8mm}p{57mm}X@{}}
\toprule
№ & Ответ & Контрольная идея \\
\midrule
1 & $x=\dfrac32$ & Основание $3$: $2x-2=x-\frac12$. \\
2 & $k=-2$, $b=3$; $f(-3)=9$; пересечение $(0;3)$. & $f(x)=-2x+3$. \\
3 & $\overrightarrow{AB}=(6;8)$, $\overrightarrow{AC}=(8;-6)$; $|AB|=10$; $\angle BAC=90^\circ$. & Скалярное произведение равно нулю. \\
4 & $9$ км/ч. & $\dfrac{24}{v-3}-\dfrac{24}{v+3}=2$, $v>3$. \\
5 & Первый: $6$ ч; второй: $12$ ч. & $\dfrac1x+\dfrac1{x+6}=\dfrac14$, $x>0$. \\
6 & $28\%$. & Соли $2{,}4+0{,}4=2{,}8$ кг, смеси $10$ кг. \\
7 & $5000$ руб.; итоговое повышение на $8\%$. & Общий коэффициент $1{,}2\cdot0{,}9=1{,}08$. \\
8 & $10\%$ годовых. & $600k^2+120k=858$, $k=1{,}1$. \\
9 & Равные платежи: по $72$ тыс. руб., переплата $34$ тыс. руб. Дифференцированные: $77$ и $66$ тыс. руб., переплата $33$ тыс. руб. Вторая схема дешевле на $1$ тыс. руб. & Остаток после второго платежа равен нулю. \\
10 & $10$ разбиений; вероятность $\dfrac25=0{,}4$. & $\binom63/2=10$; благоприятных разбиений $4$. \\
\bottomrule
\end{tabularx}
\normalsize
\vfill
{\color{accent}\rule{\linewidth}{0.55pt}}\par
\vspace{3mm}
{\small\color{darkaccent}Лёвин Артём Александрович \quad $\cdot$ \quad эксклюзивно для Анастасии Павловой}
\end{document}